\documentclass[12pt]{article}
\usepackage[utf8]{inputenc}
\usepackage[T1]{fontenc}
\usepackage{hyperref}
\title{Karahawyana Language Resources for K-12 Curriculum Development}
\author{Compiled by Claude Sonnet 5.5 (claude-sonnet-5-5) for Andr\'as Kornai}
\date{}
\begin{document}
\maketitle

\section{Introduction}
Karahawyana (ISO 639-3 xkh; Glottocode kara1475) is not a separate language in current standards. The search results say that xkh was retired in ISO 639-3 as a duplicate of Waiwai [waw] (change request 2015-038, effective 2016-01-15). Glottolog lists Karahawyana as unattested. It is treated as a variety or dialect of Waiwai, a Cariban language of northern Brazil, with some speakers in Guyana and Suriname. Joshua Project reports a Karafawyana group in Brazil (Par\'a and Amazonas) of about 50 people, and notes that they may have assimilated into larger neighbouring groups.

\section{Primary Educational Resources}
No educational resources specific to Karahawyana were found. Any curriculum work would have to use Waiwai materials. None were identified in this search.

\section{Linguistic and Reference Materials}
\begin{itemize}
\item Glottolog entry for Karahawyana: \url{https://glottolog.org/resource/languoid/id/kara1475}
\item Wikipedia, Waiwai language: \url{https://en.wikipedia.org/wiki/Waiwai_language}
\item ISO 639-3 retirement record for xkh: \url{https://gfzpublic.gfz-potsdam.de/cone/iso639-3/resource/xkh}
\end{itemize}

\section{Religious and Community Resources}
\begin{itemize}
\item Joshua Project profile, Karafawyana in Brazil: \url{https://joshuaproject.net/people_groups/12522/BR}
\end{itemize}

\section{Digital and Audio Resources}
None found.

\section{Curriculum Development Context}
Little exists under this name. The practical route is to look for Waiwai resources and for any remaining Karahawyana-identifying community. A person with local knowledge would need to confirm whether such a community still exists.
\end{document}